% theme: quadratic_equations_number_problems
\MajorQuestionLesson{二次方程式と数の利用問題}
\SetRowSpace{6mm}

\TypeQuestionSingle{次の問いに答えなさい。}{kai}
\QQ{$x=-2$ が $2x^2+(a-4)x-2a=0$ の解であるとき、$a$ の値と、もう一方の解を求めなさい。}{$a=4$、もう一方の解は $x=2$}
\QQ{$x=3$ が $x^2+(a-5)x-2a=0$ の解であるとき、$a$ の値と、もう一方の解を求めなさい。}{$a=6$、もう一方の解は $x=-4$}
\QQ{$x=-1$ が $x^2+(a+2)x+3a-1=0$ の解であるとき、$a$ の値と、もう一方の解を求めなさい。}{$a=1$、もう一方の解は $x=-2$}
\QQ{$x^2+ax+b=0$ の2つの解が $2$ と $-3$ であるとき、$a,\ b$ の値を求めなさい。}{$a=1,\ b=-6$}
\QQ{$2x^2+ax+b=0$ の2つの解が $1$ と $4$ であるとき、$a,\ b$ の値を求めなさい。}{$a=-10,\ b=8$}

\TypeQuestionSingle{次の問いに答えなさい。}{arukazu}
\QQ{ある数に3をたしてから2乗すると64になった。この数をすべて求めなさい。}{$5,\ -11$}
\QQ{ある数の2乗は、その数の3倍より10大きい。この数をすべて求めなさい。}{$5,\ -2$}
\QQ{ある数を2乗してから7をたすべきところを、7をたしてから2倍したため、正しい値より8小さくなった。この数をすべて求めなさい。}{$5,\ -3$}

\TypeQuestionSingle{次の問いに答えなさい。}{gimmi}
\QQ{ある正の数がある。その数の2乗は、その数の7倍より18大きい。この数を求めなさい。}{$9$}
\QQ{ある正の数がある。その数を2乗すると、その数に6をたしたものの4倍に等しい。この数を求めなさい。}{$6$}
\QQ{ある正の整数がある。その数の2乗からその数をひくと90になる。この整数を求めなさい。}{$10$}

\LessonPageBreak

\SetRowSpace{6mm}
\TypeQuestionSingle{次の問いに答えなさい。}{renzoku}
%この大問は吟味が必要なものも混ぜてください。
\QQ{連続する2つの自然数があり、その積は156である。この2つの自然数を求めなさい。}{$12,\ 13$}
\QQ{連続する2つの整数があり、その積は56である。この2つの整数の組をすべて求めなさい。}{$-8,\ -7$ と $7,\ 8$}
\QQ{連続する2つの正の偶数があり、その積は168である。この2つの偶数を求めなさい。}{$12,\ 14$}
\QQ{和が15である2つの正の数があり、その積は54である。この2つの数を求めなさい。}{$6,\ 9$}
\QQ{差が5である2つの正の整数があり、その積は84である。この2つの整数を求めなさい。}{$7,\ 12$}
\QQ{一方がもう一方より2大きい2つの正の整数があり、2つの数の2乗の和は100である。この2つの整数を求めなさい。}{$6,\ 8$}

\TypeQuestionSingle{次の問いに答えなさい。}{keta}
%合計4題うち２つは２ケタ、１つは３ケタにしてください。
%問題文が同じすぎるので、ある程度変えてください。「和が〜」じゃなくて「差が〜」でもいいですし、３桁は１０の位固定でもいいですね。○倍より〜〜大きい、一辺倒ではなくここにもバリエーションがほしいです。
\QQ{百の位の数と一の位の数の和が9で、十の位の数が4である3けたの自然数がある。この自然数は、百の位の数と一の位の数の積の22倍より5大きい。この自然数を求めなさい。}{$445$}
\QQ{百の位の数と一の位の数の和が7で、十の位の数が3である3けたの自然数がある。この自然数は、百の位の数と一の位の数の積の52倍より12大きい。この自然数を求めなさい。}{$532$}
\QQ{百の位の数と一の位の数の和が10で、十の位の数が6である3けたの自然数がある。この自然数は、百の位の数と一の位の数の積の17倍より10大きい。この自然数を求めなさい。}{$367$}
\QQ{百の位の数と一の位の数の和が9で、十の位の数が5である3けたの自然数がある。この自然数は、百の位の数と一の位の数の積の55倍より18小さい。この自然数を求めなさい。}{$752$}
